\documentclass[11pt]{article}
\usepackage[margin=0.85in]{geometry}
\usepackage[T1]{fontenc}
\usepackage[utf8]{inputenc}
\usepackage{lmodern}
\usepackage{microtype}
\usepackage{longtable,booktabs,array,tabularx,enumitem,float,fancyvrb,xcolor}
\usepackage[hyphens]{url}
\usepackage[colorlinks=true,linkcolor=blue!45!black,urlcolor=blue!45!black]{hyperref}
\providecommand{\tightlist}{\setlength{\itemsep}{0pt}\setlength{\parskip}{0pt}}
\definecolor{shadecolor}{RGB}{248,248,248}
\newenvironment{Shaded}{\begin{quote}\small}{\end{quote}}
\DefineVerbatimEnvironment{Highlighting}{Verbatim}{commandchars=\\\{\}}
\newcommand{\KeywordTok}[1]{#1}
\newcommand{\DataTypeTok}[1]{#1}
\newcommand{\DecValTok}[1]{#1}
\newcommand{\NormalTok}[1]{#1}
\newcommand{\OperatorTok}[1]{#1}
\newcommand{\StringTok}[1]{#1}
\newcommand{\CommentTok}[1]{\textit{#1}}
\newcommand{\BuiltInTok}[1]{#1}
\newcommand{\ExtensionTok}[1]{#1}
\newcommand{\AttributeTok}[1]{#1}
\newcommand{\ControlFlowTok}[1]{#1}
\newcommand{\FunctionTok}[1]{#1}
\newcommand{\ImportTok}[1]{#1}
\newcommand{\InformationTok}[1]{#1}
\newcommand{\WarningTok}[1]{#1}
\newcommand{\AlertTok}[1]{#1}
\newcommand{\ErrorTok}[1]{#1}
\newcommand{\ConstantTok}[1]{#1}
\newcommand{\SpecialCharTok}[1]{#1}
\newcommand{\VerbatimStringTok}[1]{#1}
\newcommand{\VerbatimTok}[1]{#1}
\newcommand{\VariableTok}[1]{#1}
\newcommand{\OtherTok}[1]{#1}
\setlength{\LTleft}{0pt}
\setlength{\LTright}{0pt}
\setlength{\emergencystretch}{3em}
\sloppy
\title{Humanvoice: Implementation and Evidence Dossier}
\author{Generated from the humanvoice evidence audit}
\date{2026-08-24T07:18:14+00:00}
\begin{document}
\hypertarget{humanvoice-implementation-and-evidence-dossier}{%
\section{Humanvoice implementation and evidence dossier}\label{humanvoice-implementation-and-evidence-dossier}}

\textbf{Status: RELEASE BLOCKED}\\
Generated by \texttt{tools/survey\_audit.py} at \texttt{2026-08-24T07:18:14+00:00} under protocol \texttt{HV-AUDIT-2026-08-22}.

This is a backstage implementation record, not the reader-facing product proposal. It records what was searched and verified, the requirements supported by those records, the proposed evaluation, and the evidence still required before stronger claims can be made.

\hypertarget{executive-decision}{%
\subsection{Executive decision}\label{executive-decision}}

Build humanvoice as a reader-acceptance workflow with a deterministic measurement layer and a separately validated feedback layer. The system should help an author or editor find register leakage, structural repetition, defensive prose, unsupported claims, citation errors, and meaning drift. It must not rewrite automatically, issue authorship verdicts, or substitute a model score for a human acceptance decision.

The immediate funding decision is a bounded implementation and pilot, not unrestricted deployment. Release of the product as evidence of improved writing is conditional on the gates in this document and on a preregistered comparison with the current workflow.

\hypertarget{evidence-at-a-glance}{%
\subsection{Evidence at a glance}\label{evidence-at-a-glance}}

\begin{itemize}
\tightlist
\item
  \textbf{Search events:} \texttt{183}. Network queries and local/import events are recorded.
\item
  \textbf{Literature candidates:} \texttt{11733}. This is a discovery pool, not an included-study count.
\item
  \textbf{Evidence rows:} \texttt{48}; \texttt{34} marked load-bearing, including \texttt{17} empirical/review rows subject to the full-text gate; \texttt{0} load-bearing rows independently verified.
\item
  \textbf{Empirical evidence profile:} \texttt{3} systematic reviews, \texttt{4} field experiments, and \texttt{14} other primary studies; source type is not a quality rating.
\item
  \textbf{Bibliographic identity:} \texttt{0} DOI-bearing rows resolved to matching Crossref records; this does not verify findings.
\item
  \textbf{Checked-in bibliography identity:} \texttt{0/25} actual DOI rows resolved; arXiv identifiers are not counted as DOIs.
\item
  \textbf{Critical appraisal:} \texttt{0} rows completed; pending appraisal prevents release-ready evidence claims.
\item
  \textbf{Software records:} \texttt{28}. Metadata inventory; effectiveness is a separate question.
\item
  \textbf{Domain imports:} \texttt{11655} total; \texttt{11655} carry complete export provenance.
\item
  \textbf{Broad known-item recall:} \texttt{6/14} (\texttt{42.9\%}). Seed-only and exact challenge matches do not count.
\item
  \textbf{Challenge-query recall:} \texttt{0/14} (\texttt{0.0\%}). This confirms identity lookup, not search saturation.
\item
  \textbf{Release gates:} \texttt{5/15}. Status remains blocked until the remainder are independently satisfied.
\end{itemize}

\hypertarget{decision-requested}{%
\subsection{Decision requested}\label{decision-requested}}

Authorize the following reversible phase:

\begin{enumerate}
\def\labelenumi{\arabic{enumi}.}
\tightlist
\item
  Freeze the protocol, schemas, and protected-source policy in version control.
\item
  Implement the format-aware diagnostics and review packet described below.
\item
  Assemble an adjudicated economics writing corpus from the internal records and external reference material.
\item
  Run one live-document feasibility test, then a preregistered, power-justified paired or stepped rollout against the current workflow.
\item
  Stop or revise the project if the kill criteria are met; do not expand the system because a surface score improves.
\end{enumerate}

\hypertarget{why-this-product}{%
\subsection{Why this product}\label{why-this-product}}

The internal failure record shows a repeated sequence: internal workflow language leaks into exposition; mechanical surface cues obscure missing explanation; defensive qualifications replace mechanisms; and readers reject documents that automated checks accepted. The accepted SMEwallet units also show that adding length and connective reasoning can improve comprehension, so compression is not a valid quality proxy.

The external evidence changes the implementation in six ways:

\begin{itemize}
\tightlist
\item
  Controlled writing experiments find productivity and quality gains from assistance, but those outcomes are task-specific and do not establish reader acceptance.
\item
  Feedback studies find that models can produce specific comments while missing the most important problem or misjudging whether criticism is warranted.
\item
  Co-writing studies find risks of convergence, disclosure effects, and loss of agency or ownership.
\item
  Citation and review studies show that fluent output can contain unverifiable references and that model judges remain weak on long, critical evaluation.
\item
  Two systematic reviews find more consistent surface and process benefits than higher-order writing gains, and warn that validity and fairness remain unsettled for consequential evaluation.
\item
  In-the-wild interaction traces and workplace deployments show that behavior and effects vary by intent, task, and writer skill; an average effect is not an expert-writing result.
\end{itemize}

The product therefore optimizes the human decision process, not a detector score or a generic notion of human-likeness.

\hypertarget{literature-synthesis-by-decision-question}{%
\subsection{Literature synthesis by decision question}\label{literature-synthesis-by-decision-question}}

The evidence is heterogeneous, so the proposal uses it as a set of bounded design constraints rather than as a pooled effect estimate.

\begin{itemize}
\tightlist
\item
  \textbf{Does assistance improve work?} Controlled experiments report task-specific gains in speed or evaluated output quality, including professional writing tasks and workplace email. A deployment to 5,172 support agents also found heterogeneous effects, with larger gains for less-experienced workers and small quality losses among the most skilled. Those estimates do not identify the value of this repository's reader-acceptance workflow, so they motivate a skill-stratified pilot rather than an ROI claim.
\item
  \textbf{Can a model give useful revision feedback?} Recent writing-feedback studies show that models can produce specific comments, but they can miss the most important defect, vary by writer proficiency, and disagree with human judgments. Systematic reviews of 25 and 96 empirical studies find more consistent surface/process gains than higher-order writing gains. The product must rank issues, expose uncertainty, and measure the next revision against a domain rubric.
\item
  \textbf{What can be lost?} Co-writing and disclosure studies point to agency, ownership, reader-perception, and collective-diversity effects. These are not cosmetic risks: a workflow that makes documents more uniform while reducing reader trust is a failed intervention.
\item
  \textbf{Can a machine judge replace a reader?} Citation-reliability and long-form review studies make that substitution unsafe. Model judges may screen, but acceptance remains a named human decision and every load-bearing claim needs an evidence anchor.
\item
  \textbf{What process should be measured?} In-the-wild writing traces show recurring intent revision, exploration, questioning, style adjustment, and content injection. The workflow must retain suggestion, acceptance, edit, rejection, and rollback events; output-only comparisons cannot explain how quality or burden changed.
\item
  \textbf{How should the survey be trusted?} PRISMA-S, OpenAlex, and Crossref provide reporting and infrastructure patterns for replayable searches; they do not make public API coverage equivalent to licensed economics or bibliographic databases. The separate domain-import and specialist-coverage gates are therefore substantive, not paperwork.
\end{itemize}

\hypertarget{review-method-and-current-certainty}{%
\subsection{Review method and current certainty}\label{review-method-and-current-certainty}}

The protocol freezes inclusion classes, seven exclusion codes, dual-review decisions, design-specific critical-appraisal tools, and a known-item recall challenge. Discovery records remain distinct from included evidence. The first-pass queue is deterministic and bounded; it is a work allocation device, not a claim that screening is complete.

DOI resolution verifies title identity only. Empirical findings become release-ready only after full-text inspection, a source anchor for the extracted claim, design-specific appraisal, limitations, and independent verification. This separation prevents a valid citation from being mistaken for a valid inference.

\hypertarget{search-coverage-and-failure-accounting}{%
\subsection{Search coverage and failure accounting}\label{search-coverage-and-failure-accounting}}

A large candidate count is not evidence of exhaustive retrieval. The run records every source event and its exact status so an API quota, an authentication limit, or a missing licensed database cannot be mistaken for a negative result.

\begin{itemize}
\tightlist
\item
  \textbf{crossref:} 0/46 events completed retrieval or import; status ledger: network-error=46.
\item
  \textbf{crossref-bibliography:} 0/25 events completed retrieval or import; status ledger: network-error=25.
\item
  \textbf{crossref-doi:} 0/16 events completed retrieval or import; status ledger: network-error=16.
\item
  \textbf{github:} 0/14 events completed retrieval or import; status ledger: network-error=14.
\item
  \textbf{iclr-proceedings:} 0/1 events completed retrieval or import; status ledger: network-error=1.
\item
  \textbf{manual-import:} 6/6 events completed retrieval or import; status ledger: imported=6.
\item
  \textbf{nber:} 0/1 events completed retrieval or import; status ledger: network-error=1.
\item
  \textbf{openalex:} 0/46 events completed retrieval or import; status ledger: network-error=46.
\item
  \textbf{package-registry:} 0/28 events completed retrieval or import; status ledger: network-error=28.
\end{itemize}

OpenAlex rate-limit responses, when present, are preserved as raw JSON. Broad recall is computed across all declared discovery sources and must be read alongside this status ledger. The ACL, NBER, and ICLR public exports improve domain coverage but are explicitly reported as filtered public slices, not substitutes for licensed indexes.

\hypertarget{product-boundary}{%
\subsection{Product boundary}\label{product-boundary}}

\hypertarget{in-scope}{%
\subsubsection{In scope}\label{in-scope}}

\begin{itemize}
\tightlist
\item
  A protected-source representation for LaTeX, Markdown, citations, equations, tables, code, and quoted material.
\item
  Deterministic diagnostics for register leakage, cadence, defensive language, unsupported specificity, citation identity, compilation, and version drift.
\item
  A review packet containing quoted spans, context, evidence links, and an editorial question. Diagnostics flag; a human or explicitly invoked editor decides.
\item
  A reader-facing acceptance workflow with low-context review, blinded pairwise comparison, and a recorded human verdict.
\item
  Corpus-level monitoring of convergence and machine-typical language, without per-document authorship classification.
\end{itemize}

\hypertarget{out-of-scope}{%
\subsubsection{Out of scope}\label{out-of-scope}}

\begin{itemize}
\tightlist
\item
  Automatic authorship or misconduct judgments.
\item
  Detector evasion, watermark removal, or optimization against a vendor's classifier.
\item
  Autonomous generation of claims, citations, data, or references.
\item
  Uploading confidential manuscripts to an unapproved remote provider.
\item
  A claim that any tool reduces human-feedback rounds before the pilot measures it.
\end{itemize}

\hypertarget{operating-model}{%
\subsection{Operating model}\label{operating-model}}

The workflow has four actors. The author owns claims, evidence, and final wording. The diagnostic layer measures observable properties and emits spans. The editor or writing agent proposes bounded changes while preserving protected content. The reader is the acceptance authority. A research steward maintains the corpus, protocol, package pins, and disclosure record.

Every intervention produces two artifacts: a reader-facing draft and an internal provenance record. The reader-facing draft should not contain row IDs, gate names, or process vocabulary. The provenance record stores the source hash, tool versions, prompts or rule sets, changed spans, reviewer decisions, and disclosure status.

\hypertarget{system-architecture}{%
\subsection{System architecture}\label{system-architecture}}

\begin{verbatim}
source files -> format parser -> protected document model
                              |
                              +--> deterministic diagnostics
                              |       (flags and measurements)
                              +--> evidence/citation resolver
                              |
                              +--> meaning-preserving diff
                              |
                              +--> review packet -> human decision
                                                      |
                         corpus + verdicts <-----------+
                              |
                         held-out evaluation and drift monitoring
\end{verbatim}

The measurement layer is intentionally modular. A package can be removed without changing the acceptance protocol. A model can be swapped without changing the protected-source or evidence schemas. No diagnostic is allowed to promote a document by itself.

\hypertarget{functional-requirements}{%
\subsection{Functional requirements}\label{functional-requirements}}

Each requirement has an acceptance test and an evidence trail. The IDs are stable so implementation issues and pilot results can refer back to them.

\begin{itemize}
\tightlist
\item
  \textbf{R1 Reader-facing separation.} Keep the reader-facing manuscript separate from the audit and provenance record, linked by immutable document and version identifiers; internal workflow vocabulary must not leak into the manuscript. \textbf{Acceptance:} A low-context reader can state the purpose, object, and requested action, the leak fixture finds no forbidden internal register, and every released file traces to a source hash and version. \textbf{Evidence:} E-INT-001, E-INT-002. \textbf{Phase:} P1,P4.
\item
  \textbf{R2 Configurable policy precedence.} Resolve conflicts in the order correctness, source fidelity, domain meaning, reader comprehension, then template regularity; store project and format overrides in configuration rather than manuscript prose. \textbf{Acceptance:} Conflict fixtures resolve in the declared order and LaTeX-specific configuration disables rules that would alter protected syntax. \textbf{Evidence:} E-INT-002, E-SW-003, E-SW-013. \textbf{Phase:} P1.
\item
  \textbf{R3 Located diagnostics without automatic rewriting.} Diagnostics emit categories, source spans, measurements, context, and editorial questions; only an explicit author or editor action may change text. \textbf{Acceptance:} Every diagnostic fixture returns stable locations and no command mutates its input; accepted edits are remeasured and recorded separately. \textbf{Evidence:} E-INT-002, E-SW-003, E-SW-013. \textbf{Phase:} P2,P4.
\item
  \textbf{R4 Layered diagnosis and repair.} Diagnose and repair register leakage, mechanical repetition, defensive prose, substance, and reader comprehension in that order while retaining intermediate results. \textbf{Acceptance:} Historical replay preserves a result for every layer and prevents a later-layer score from concealing an unresolved earlier defect. \textbf{Evidence:} E-INT-001, E-EXT-007. \textbf{Phase:} P2,P4.
\item
  \textbf{R5 Protected meaning-preserving diff.} Represent equations, labels, citations, numbers, claims, code, tables, and quoted text as protected content and explain every substantive change between versions. \textbf{Acceptance:} Repair replay produces zero unexplained protected-token changes and blocks promotion on any unexplained deletion, insertion, or mutation. \textbf{Evidence:} E-INT-001, E-SW-006, E-SW-007, E-SW-008. \textbf{Phase:} P1,P2.
\item
  \textbf{R6 Citation identity and claim support.} Resolve bibliographic identity separately from source-to-claim support, retain a sentence-level claim anchor, and block release on an unresolved load-bearing citation. \textbf{Acceptance:} Seeded valid, nonexistent, wrong-DOI, duplicate, and title-mismatch cases are classified correctly, and every load-bearing claim has a human-verified full-text anchor. \textbf{Evidence:} E-EXT-015, E-METHOD-005, E-METHOD-007. \textbf{Phase:} P0,P2.
\item
  \textbf{R7 Calibrated model judging.} Use model judges only for construct-specific screening after pair-order randomization, length control, cross-family checking, and calibration against recorded human labels. \textbf{Acceptance:} A frozen human-labelled set reports agreement and order, verbosity, and family-sensitivity diagnostics before a judge result can enter analysis. \textbf{Evidence:} E-EXT-008, E-EXT-016, E-EXT-018. \textbf{Phase:} P3.
\item
  \textbf{R8 Named human acceptance.} A human reader, not a surface metric or model judge, is the promotion authority for a specific document version. \textbf{Acceptance:} Every promoted document records the pseudonymous reader, immutable document hash, rubric response, decision, requested changes, and elapsed time. \textbf{Evidence:} E-INT-001, E-METHOD-006. \textbf{Phase:} P3,P4,P5.
\item
  \textbf{R9 Held-out behavioral evaluation.} Evaluate diagnostic packages and the complete workflow on documents not used for rule tuning, and separate surface, process, higher-order, and acceptance outcomes. \textbf{Acceptance:} Tuning and holdout hashes are frozen before execution, each package has a human-labelled scorecard, and no smoke test is reported as effectiveness. \textbf{Evidence:} E-INT-002, E-EXT-017, E-EXT-018. \textbf{Phase:} P0,P3,P5.
\item
  \textbf{R10 Minimal process surface.} Use one bounded plan and one result record per work package; do not expose audit ledgers, gate names, or orchestration state in reader-facing prose. \textbf{Acceptance:} The release packet contains the required provenance records while the manuscript passes the internal-register leak test and a reader can identify the argument without process context. \textbf{Evidence:} E-INT-001, E-INT-002. \textbf{Phase:} P1,P4.
\item
  \textbf{R11 Privacy and disclosure.} Keep protected manuscripts local by default and record provider, model, purpose, retention policy, disclosure status, and approval for every remote call. \textbf{Acceptance:} An unapproved remote route is blocked, approved calls produce a complete disclosure record, and the minimum-content retention policy is enforced. \textbf{Evidence:} E-METHOD-004, E-METHOD-006. \textbf{Phase:} P1,P4,P5.
\item
  \textbf{R12 Population-level convergence without attribution.} Monitor aggregate vocabulary and structural convergence against an appropriate pre-assistance baseline, but never classify the authorship of an individual document. \textbf{Acceptance:} The drift command refuses single-document analysis, suppresses small cells, and reports uncertainty and baseline sensitivity for aggregate trends. \textbf{Evidence:} E-EXT-004, E-EXT-005, E-EXT-012. \textbf{Phase:} P2,P3,P5.
\item
  \textbf{R13 Interaction telemetry.} Record writer intent, suggestion, acceptance, edit, rejection, and rollback events without retaining unnecessary manuscript content. \textbf{Acceptance:} End-to-end replay reconstructs the intervention path from hashes and event records, distinguishes edited acceptance from direct acceptance, and applies the retention class. \textbf{Evidence:} E-EXT-002, E-EXT-019, E-EXT-021. \textbf{Phase:} P3,P4,P5.
\item
  \textbf{R14 Heterogeneous outcomes.} Report effects by baseline writer skill, document task and genre, and reader role; do not deploy from an average effect that conceals subgroup harm. \textbf{Acceptance:} The preregistered analysis reports uncertainty for each prespecified stratum and applies the stop rule if expert writers or a task class are harmed. \textbf{Evidence:} E-EXT-013, E-EXT-017, E-EXT-018, E-EXT-020. \textbf{Phase:} P3,P5.
\item
  \textbf{R15 Format-aware parser contract.} All parser adapters must emit the same source-map and protected-span contract; no parser is authoritative merely because it accepts a document. \textbf{Acceptance:} Pandoc, pylatexenc, plasTeX, LaTeXML, and tree-sitter candidates are benchmarked on malformed, unknown-macro, nested-environment, comment, math, citation, location, and round-trip fixtures before selection. \textbf{Evidence:} E-SW-006, E-SW-007, E-SW-014, E-SW-015, E-SW-016. \textbf{Phase:} P1.
\item
  \textbf{R16 Evidence-tiered package adoption.} Treat package availability, operational suitability, and writing effectiveness as separate decisions; only a held-out human-labelled benchmark supports adoption. \textbf{Acceptance:} Every adoption candidate has one provenance record and one explicit hold or adopt decision, and adoption is impossible unless its benchmark status is effectiveness-pass and independent review is recorded. \textbf{Evidence:} E-INT-002, E-SW-003, E-SW-013. \textbf{Phase:} P0,P1,P3.
\end{itemize}

\hypertarget{implementation-contracts}{%
\subsection{Implementation contracts}\label{implementation-contracts}}

The core is a versioned document and evidence model, not a chain of prompts. Every command reads immutable source hashes, emits schema-versioned JSON, and records the tool, configuration, and model version that produced each result. Style findings return data rather than silently changing a file; only invariant, build, privacy, or unresolved load-bearing citation failures may return a blocking status.

The minimum records are:

\begin{itemize}
\tightlist
\item
  \textbf{DocumentRecord:} document and version IDs, source hash, format, source map, protected spans, policy/configuration version, and disclosure class.
\item
  \textbf{DiagnosticFinding:} stable finding ID, rule and category, source span, excerpt hash, surrounding context, measurement, editorial question, and disposition. Severity is operational priority, not a quality score.
\item
  \textbf{ChangeRecord:} before/after hashes and spans, editor or tool identity, stated reason, protected-content impact, and author disposition.
\item
  \textbf{CitationRecord:} citation key, supplied and canonical DOI/URL metadata, identity status, sentence-level claim anchors, full-text verification, and reviewer.
\item
  \textbf{ReviewDecision:} pseudonymous reader, blinded condition where feasible, document hash, rubric responses, acceptance or revision decision, confidence, requested changes, and elapsed time.
\item
  \textbf{InteractionEvent:} session and document version, actor, intent class, suggestion hash, accept/edit/reject/rollback event, editing distance, timestamp, and retention class. Store the minimum manuscript content needed for audit.
\end{itemize}

The seven command contracts are:

\begin{itemize}
\tightlist
\item
  \textbf{\texttt{hv\ leak}:} source plus project policy to located internal-register and disclosure findings. It must replay the preserved rejected/accepted cases and report locations without editing.
\item
  \textbf{\texttt{hv\ rhythm}:} source and rendered prose to sentence, paragraph, opener, transition, and convergence measurements. Thresholds are genre-configured and never certify quality.
\item
  \textbf{\texttt{hv\ defensive}:} source to contextualized candidate spans and editorial questions. Precision, recall, and reviewer burden are measured on labelled tuning and untouched holdout sets before activation.
\item
  \textbf{\texttt{hv\ diff}:} two DocumentRecords to protected-token censuses, changed spans, and an explain-every-removal ledger. Any unexplained equation, number, label, citation, or claim change blocks promotion.
\item
  \textbf{\texttt{hv\ cite}:} bibliography plus claim contexts to canonical identity results and a human verification queue. Seeded valid, nonexistent, wrong-DOI, duplicate, and title-mismatch fixtures must be classified correctly.
\item
  \textbf{\texttt{hv\ gate}:} source to reproducible build status, unresolved-reference report, layout warnings, PDF hash, and selected rendered-page images. A PDF that compiles but fails visual inspection does not pass.
\item
  \textbf{\texttt{hv\ drift}:} a sufficiently sized document population to aggregate vocabulary, structure, and convergence trends. It refuses single-document attribution and suppresses small cells.
\end{itemize}

Parser adapters implement the same source-map contract. Pandoc and pylatexenc are the first benchmark pair; plasTeX is the macro-expanding fallback, LaTeXML the heavier conversion fallback, and tree-sitter-latex the incremental-editor candidate. No adapter becomes authoritative until malformed, unknown-macro, nested-environment, comment, math, citation, and round-trip fixtures pass.

\hypertarget{evaluation-design}{%
\subsection{Evaluation design}\label{evaluation-design}}

The primary outcome is \textbf{human-feedback rounds to reader acceptance}. Secondary outcomes are reader reconstruction accuracy, substantive revision rate, meaning-preservation violations, time to accepted unit, false-positive burden, suggestion acceptance and editing cost, cross-document convergence. The pilot must compare the complete workflow with a documented baseline, use the same document task distribution, blind the reader to condition where feasible, and report paired differences with uncertainty rather than a single average.

One live document is a feasibility test only: one live document for integration, telemetry, privacy, stop-rule, and burden testing only. The efficacy phase uses paired or stepped rollout over live documents, selected before outcomes are inspected. Its sample size is determined by this rule: simulate power or precision from the observed baseline distribution of feedback rounds and reader clustering; do not choose an arbitrary document count. Prespecified strata are baseline writer skill, document task and genre, reader role.

The corpus should contain at least four strata: accepted human-authored economics prose, rejected and repaired internal drafts, AI-assisted drafts with interaction traces, and deliberately corrupted or adversarial examples. Each item needs a provenance record, protected-source map, defect labels, and reader verdicts. Tuning and held-out evaluation splits must be frozen before benchmark results are reported.

The reader study should record: reconstruction of purpose and contribution, perceived substance, factual/citation trust, willingness to act, revision requests, confidence, and time. A second reader sample estimates inter-reader agreement. Model judges can reduce screening cost only after their position, verbosity, self-preference, and domain biases are measured on the human-labelled subset.

The causal question is narrow: does the workflow reduce the number of human-feedback rounds required to reach acceptance without increasing meaning violations or reader time? It is not whether a document receives a higher detector score.

\hypertarget{software-strategy}{%
\subsection{Software strategy}\label{software-strategy}}

The survey separates three claims that are often conflated: availability, operational suitability, and effectiveness. Registry and repository metadata support only the first. A pinned, reproducible adapter supports the second. A held-out, human-labelled benchmark supports the third.

The adoption candidates are listed below. A package remains survey-only until the stated benchmark is complete.

\begin{itemize}
\tightlist
\item
  \textbf{sloptrim} (ai-tell-diagnostic): benchmark status \texttt{smoke-pass}. Before adoption, test protected-region behavior, false positives, runtime, license, and held-out effectiveness.
\item
  \textbf{Vale} (prose-linter): benchmark status \texttt{not-run}. Before adoption, test protected-region behavior, false positives, runtime, license, and held-out effectiveness.
\item
  \textbf{textlint} (prose-linter): benchmark status \texttt{not-run}. Before adoption, test protected-region behavior, false positives, runtime, license, and held-out effectiveness.
\item
  \textbf{LanguageTool} (grammar-style): benchmark status \texttt{not-run}. Before adoption, test protected-region behavior, false positives, runtime, license, and held-out effectiveness.
\item
  \textbf{Textidote} (latex-writing): benchmark status \texttt{not-run}. Before adoption, test protected-region behavior, false positives, runtime, license, and held-out effectiveness.
\item
  \textbf{LTeX+ LS} (latex-writing): benchmark status \texttt{not-run}. Before adoption, test protected-region behavior, false positives, runtime, license, and held-out effectiveness.
\item
  \textbf{Pandoc} (document-conversion): benchmark status \texttt{not-run}. Before adoption, test protected-region behavior, false positives, runtime, license, and held-out effectiveness.
\item
  \textbf{pylatexenc} (latex-parsing): benchmark status \texttt{not-run}. Before adoption, test protected-region behavior, false positives, runtime, license, and held-out effectiveness.
\item
  \textbf{latexdiff} (document-diff): benchmark status \texttt{not-run}. Before adoption, test protected-region behavior, false positives, runtime, license, and held-out effectiveness.
\item
  \textbf{humanvoice diagnostics} (local product code): build status. Acceptance requires replay fixtures, reader agreement, and the pilot endpoint.
\end{itemize}

The initial stack should prefer local, format-aware, composable components. Pandoc and pylatexenc compete in the first protected-source parser benchmark; plasTeX, LaTeXML, and tree-sitter-latex are explicit fallbacks for macro expansion, structured conversion, and incremental editor spans. latexdiff is a visual-review supplement, never the invariant checker. Vale-style rules, textlint, LanguageTool, TeXtidote, and the maintained LTeX+ LS remain replaceable diagnostic adapters; archived ltex-ls and the deep review's rejected proselint dependency are not adoption candidates. GROBID is optional PDF ingestion, textstat is descriptive only, Inspect AI is an optional model-evaluation harness, and llama.cpp is a privacy-capable runtime whose model quality must be evaluated separately. Remote model calls are optional and must be explicit. AI detectors remain survey-only unless a future population-level validation supports a narrowly defined monitoring use.

The machine-readable \texttt{package\_selection.json} records one canonical decision per adoption candidate. The current policy is hold-pending-held-out-benchmark; a smoke pass is not an adoption decision.

\hypertarget{implementation-plan}{%
\subsection{Implementation plan}\label{implementation-plan}}

\begin{itemize}
\tightlist
\item
  \textbf{P0, weeks 1-4, evidence completion:} specialist and licensed domain exports, dual screening and adjudication, full-text claim anchors, design-specific critical appraisal, and package selection record. Exit when the evidence gates and provenance checks pass.
\item
  \textbf{P1, weeks 1-3, policy and protected source:} schemas, disclosure policy, parser adapter benchmark, protected-region model, and baseline map. Exit when round-trip and malformed-source fixtures pass.
\item
  \textbf{P2, weeks 2-6, deterministic diagnostics:} implement \texttt{leak}, \texttt{rhythm}, \texttt{defensive}, \texttt{cite}, \texttt{diff}, \texttt{gate}, and \texttt{drift} contracts. Exit when historical replay, invariant, and citation fixtures pass.
\item
  \textbf{P3, weeks 3-7, corpus and harness:} labelled repair pairs, reader rubric, and blinded pairwise harness. Exit when inter-rater and holdout splits are recorded.
\item
  \textbf{P4, weeks 6-9, workflow integration:} review packet, editor interface, provenance record, and CI entry point. Exit when one document runs end to end.
\item
  \textbf{P5, weeks 9-12, feasibility:} run one live document end to end and report integration failures, privacy events, telemetry completeness, and burden; do not estimate efficacy from it. Then preregister the powered paired or stepped rollout. The efficacy schedule follows the power/precision analysis and live-document arrival rate rather than an arbitrary 12-week promise.
\end{itemize}

Resource planning assumption: one product engineer, one research/evaluation lead, and part-time domain reviewer; planning assumption, not an observed cost. This is a budgeting placeholder, not a measured cost estimate. Planning boundary: Twelve weeks covers implementation and one feasibility document; efficacy timing follows the power or precision analysis and live-document arrival rate.

\hypertarget{economic-decision-model}{%
\subsection{Economic decision model}\label{economic-decision-model}}

The investment case should be computed from observed workflow data. Let \texttt{r0} and \texttt{r1} be baseline and treatment feedback rounds, \texttt{t0} and \texttt{t1} reader hours per round, \texttt{c\_r} the loaded cost of reader time, and \texttt{C} the one-time build and maintenance cost. The pilot supports deployment only if the expected discounted value of \texttt{(r0*t0\ -\ r1*t1)*c\_r} plus any measured quality or time benefit exceeds \texttt{C}, subject to the meaning-preservation, privacy, and disclosure constraints. No current literature result supplies these project-specific values.

This formulation prevents productivity studies from being misused as an ROI claim for humanvoice. External experiments establish that assistance can change task time or quality under particular conditions; they do not establish the value of this workflow for this repository's documents.

\hypertarget{governance-and-risks}{%
\subsection{Governance and risks}\label{governance-and-risks}}

\begin{itemize}
\tightlist
\item
  \textbf{House-style convergence:} rising cross-document similarity or loss of distinctive arguments. Control with population-level drift reports and no single-style optimization. Pause if diversity falls without reader benefit.
\item
  \textbf{Meaning drift:} changed numbers, citations, equations, or scope after editing. Control with protected diffs and human source checks. Stop on any unexplained load-bearing change.
\item
  \textbf{False-positive fatigue:} readers ignore or disable diagnostics. Measure review burden and precision by category; stop if flags cost more time than they save.
\item
  \textbf{Fluent factual error:} unsupported or mismatched references. Use DOI/URL resolution, evidence anchors, and human verification; stop on unresolved load-bearing citations.
\item
  \textbf{Judge bias:} order swaps or length padding change verdicts. Randomize pair order, control length, and calibrate on human labels; pause if the judge misses its threshold.
\item
  \textbf{Privacy or disclosure failure:} confidential text leaves the approved boundary. Use local-first routing and provider logs; stop on any unapproved transmission.
\item
  \textbf{Process capture:} the proposal becomes a ledger instead of a product. Keep reader and audit artifacts separate and use short result notes; pause if a reader cannot explain the purpose.
\end{itemize}

\hypertarget{release-decision}{%
\subsection{Release decision}\label{release-decision}}

The current run is not a release. The exact gate results are reproduced below and remain authoritative:

The gate list is intentionally short and auditable:

\begin{itemize}
\tightlist
\item
  \textbf{protocol-frozen:} \texttt{pass}. Versioned JSON protocol is checked in.
\item
  \textbf{minimum-query-count:} \texttt{pass}. 32 OpenAlex query families are specified.
\item
  \textbf{known-item-broad-recall:} \texttt{fail}. Search retrieved 6/14 known items; offline=False.
\item
  \textbf{known-item-challenge-recall:} \texttt{fail}. Challenge queries retrieved 0/14 known items; offline=False.
\item
  \textbf{manual-independent-screening:} \texttt{fail}. Two independent include/exclude/uncertain decisions and include/exclude adjudication are required for every selected queue row; excluded records also require a frozen exclusion code.
\item
  \textbf{domain-database-imports:} \texttt{pass}. At least one provenance-complete domain export is required; records=11655, provenance-complete=11655.
\item
  \textbf{specialist-domain-coverage:} \texttt{fail}. At least one provenance-complete specialist/licensed export is required (EconLit, RePEc/IDEAS, SSRN, ACM Digital Library, IEEE Xplore, Web of Science, Scopus); matched=none. Public ACL/NBER/ICLR slices do not satisfy this gate.
\item
  \textbf{full-text-load-bearing-evidence:} \texttt{fail}. Every load-bearing empirical/review record needs full-text inspection and independent verification; verified=0/17.
\item
  \textbf{critical-appraisal:} \texttt{fail}. Every load-bearing empirical/review record needs a low/some-concern design-specific appraisal, claim anchor, and limitations; acceptable=0/17.
\item
  \textbf{bibliographic-identity:} \texttt{fail}. Every DOI-bearing load-bearing evidence row must resolve to a matching Crossref title; verified=0/14.
\item
  \textbf{bibliography-doi-identity:} \texttt{fail}. Every actual DOI in the checked-in bibliography must resolve to a matching Crossref title; verified=0/25, mismatches=0, unresolved=25. arXiv identifiers are tracked separately.
\item
  \textbf{package-provenance:} \texttt{pass}. Each package must have a URL and a pinned local commit or an explicit remote-only status.
\item
  \textbf{package-held-out-benchmark:} \texttt{fail}. Every adoption candidate needs an adapter and effectiveness results on the locked corpus; smoke tests do not establish effectiveness.
\item
  \textbf{external-review:} \texttt{fail}. The protocol requires an independent librarian/domain reviewer sign-off file with status=approved.
\item
  \textbf{artifact-integrity:} \texttt{pass}. Raw-response hashes, record identities, evidence provenance, and package provenance must remain internally consistent.
\end{itemize}

A release-ready proposal requires all of the following: broad search recall at the frozen target; genuine domain exports with source metadata; two independent screeners and adjudication; full-text anchors, acceptable critical appraisal, and independent verification for every load-bearing empirical claim; held-out package effectiveness results; and an independent librarian plus economics/HCI sign-off. Passing the automated parts cannot substitute for those human inputs.

\hypertarget{reproducibility-and-handoff}{%
\subsection{Reproducibility and handoff}\label{reproducibility-and-handoff}}

Run the following from the repository root:

\begin{Shaded}
\begin{Highlighting}[]
\ExtensionTok{python3}\NormalTok{ tools/survey\_audit.py run {-}{-}online {-}{-}execute{-}tools {-}{-}build{-}pdf {-}{-}replace{-}output}
\ExtensionTok{python3}\NormalTok{ tools/survey\_audit.py validate {-}{-}run docs/survey/audit/latest}
\ExtensionTok{python3}\NormalTok{ tools/survey\_audit.py publish {-}{-}run docs/survey/audit/latest}
\end{Highlighting}
\end{Shaded}

The run folder contains the frozen protocol, raw responses and hashes, candidate registry, complete screening ledger, deterministic screening queue, review flow, evidence and DOI-identity matrices, critical-appraisal queue, software inventory, package selection record, evaluation plan, benchmark results, corpus manifest, PDF build record, artifact manifest, requirements traceability matrix, and gate results. \texttt{publish} leaves the dossier and audit report in this run folder and copies only machine-readable requirements and status summaries; it never writes a second reader-facing document or the authored product proposal.

\hypertarget{evidence-register}{%
\subsection{Evidence register}\label{evidence-register}}

The full machine-readable evidence matrix remains in \texttt{evidence.csv}; the compact register below keeps the dossier navigable while retaining the source link.

\begin{itemize}
\tightlist
\item
  \textbf{E-INT-001} (internal-case-record; full local source; questions Q1, Q2, Q3, Q6): The internal cases describe layered failures in register, mechanical tells, defensive prose, and substance. The ZLB rescue required five human-feedback rounds; SMEwallet records unit-level acceptance but explicitly does not claim whole-document acceptance. \href{../part1_failures.tex}{source}.
\item
  \textbf{E-INT-002} (internal-package-audit; full local source and checked-in vendor tree; questions Q5, Q6): The existing audit covers a narrow set of prose and AI-tell tools, records a useful sloptrim calibration, and does not establish that any tool reduces human-feedback rounds. \href{../part4_oss.tex}{source}.
\item
  \textbf{E-EXT-001} (primary-study; primary abstract and metadata; questions Q2, Q3): A preregistered online experiment assigned occupation-specific writing tasks to 453 college-educated professionals. The abstract reports lower completion time and higher evaluated output quality with ChatGPT assistance. \href{https://ui.adsabs.harvard.edu/abs/2023Sci...381..187N/abstract}{source}.
\item
  \textbf{E-EXT-002} (primary-study; primary abstract and metadata; questions Q1, Q4): CoAuthor contains interaction traces from 63 writers and 1,445 writing sessions with four GPT-3 instances, enabling analysis of language, ideation, and collaboration behavior. \href{https://arxiv.org/abs/2201.06796}{source}.
\item
  \textbf{E-EXT-003} (primary-study; primary abstract, metadata, and proceedings record; questions Q2, Q3): The paper studies generated writing feedback and evaluates it on an economic essay assignment, while explicitly noting that the effectiveness of model feedback on student revision is not settled. \href{https://aclanthology.org/2024.emnlp-main.928/}{source}.
\item
  \textbf{E-EXT-004} (primary-study; primary landing page and abstract-level record; questions Q4): The study experimentally examines individual creative performance and collective diversity under generative AI assistance. \href{https://doi.org/10.1126/sciadv.adn5290}{source}.
\item
  \textbf{E-EXT-005} (primary-study; primary article abstract and metadata; questions Q1, Q3): The study compares 50 real abstracts with 50 ChatGPT-generated abstracts and includes blinded human review alongside detector measurements. \href{https://www.nature.com/articles/s41746-023-00819-6}{source}.
\item
  \textbf{E-EXT-006} (systematic-review; primary preprint abstract and methods summary; questions Q4, Q6): The review reports a PRISMA-based analysis of 109 HCI papers selected from more than 1,600 records, with attention to agency and ownership across planning, translation, reviewing, and monitoring. \href{https://arxiv.org/abs/2504.12488}{source}.
\item
  \textbf{E-METHOD-001} (reporting-guideline; official PRISMA-S page; questions Q6): PRISMA-S specifies sixteen items for transparently reporting literature searches, including sources, strategies, dates, and record management. \href{https://www.prisma-statement.org/prisma-search}{source}.
\item
  \textbf{E-SW-001} (software-documentation; official documentation; questions Q5): Vale provides code-like linting for prose and supports configurable rules and targeted text segments. \href{https://docs.vale.sh/}{source}.
\item
  \textbf{E-SW-002} (software-documentation; official repository; questions Q5): textlint is a pluggable natural-language linter with machine-readable formatter options. \href{https://github.com/textlint/textlint}{source}.
\item
  \textbf{E-METHOD-002} (search-infrastructure; official API documentation; questions Q6): OpenAlex documents reproducible Boolean, phrase, proximity, fuzzy, and semantic search modes over scholarly works. \href{https://help.openalex.org/api/searching/}{source}.
\item
  \textbf{E-METHOD-003} (search-infrastructure; official API documentation; questions Q6): Crossref exposes searchable scholarly metadata, DOI lookup, filters, and JSON responses through a public REST API. \href{https://www.crossref.org/documentation/retrieve-metadata/rest-api/}{source}.
\item
  \textbf{E-EXT-007} (primary-study; primary abstract, metadata, and ACL proceedings record; questions Q2, Q3): The authors introduce a 1,300-story test set with intentionally corrupted writing and evaluate model feedback with automatic and human measures. Models often give specific and mostly accurate comments but fail to identify the most important problem and to choose appropriately between critical and positive feedback. \href{https://aclanthology.org/2025.acl-long.1254/}{source}.
\item
  \textbf{E-EXT-008} (primary-study; primary abstract and proceedings metadata; questions Q2, Q3): GPT-4 showed potential for assessing revision quality when detailed rubrics described common revision patterns, but agreement with human ratings varied across writing-proficiency levels and chain-of-thought prompting did not materially improve the result. \href{https://aclanthology.org/2024.bea-1.30/}{source}.
\item
  \textbf{E-EXT-009} (primary-study; primary abstract and proceedings record; questions Q2, Q3): The study compares ChatGPT feedback with volunteer educator feedback on English-proficiency essays and reports that personalized feedback remains difficult; the authors argue that the two sources can complement one another. \href{https://aclanthology.org/2024.lrec-main.144/}{source}.
\item
  \textbf{E-EXT-010} (primary-study; primary abstract, metadata, and proceedings record; questions Q4, Q6): In experiments on argumentative essays and creative stories, disclosure of AI assistance, especially content generation, lowered average quality ratings and increased variation across readers; confidence and prior AI familiarity moderated the effect. \href{https://aclanthology.org/2024.emnlp-main.279/}{source}.
\item
  \textbf{E-EXT-011} (field-experiment; NBER working-paper abstract and methods summary; questions Q2, Q4): A randomized field experiment across 66 firms and 7,137 knowledge workers found that treated users who adopted the tool spent about two fewer hours per week on email, while the study did not detect changes in the quantity or composition of tasks. \href{https://www.nber.org/papers/w33795}{source}.
\item
  \textbf{E-EXT-012} (field-experiment; ICLR proceedings abstract and metadata; questions Q2, Q4): A controlled experiment comparing unaided writing, GPT-3 assistance, and InstructGPT assistance found a statistically significant reduction in lexical and content diversity with InstructGPT, attributable mainly to model-contributed text. \href{https://proceedings.iclr.cc/paper_files/paper/2024/hash/02dec8877fb7c6aa9a79f81661baca7c-Abstract-Conference.html}{source}.
\item
  \textbf{E-EXT-013} (field-experiment; open-access published article abstract and experiment summary; questions Q2, Q3): In a preregistered field experiment with 758 consultants, GPT-4 improved speed and quality on tasks inside its capability frontier but reduced correctness on tasks outside that frontier despite higher-rated recommendations. \href{https://www.hbs.edu/ris/Publication\%20Files/dell-acqua-et-al-2026-navigating-the-jagged-technological-frontier_5c589c8c-fbb5-458f-b285-c944746cd717.pdf}{source}.
\item
  \textbf{E-EXT-014} (primary-study; open-access abstract and study summary; questions Q1, Q3): Across experiments with novice and experienced teachers, participants generally could not identify ChatGPT essays among student writing and were overconfident in their judgments; source disclosure also affected quality judgments heterogeneously. \href{https://www.sciencedirect.com/science/article/pii/S2666920X24000109}{source}.
\item
  \textbf{E-EXT-015} (primary-study; open-access article abstract and methods/results summary; questions Q1, Q3, Q6): Two researchers prompted GPT-3.5 to write introductions with citations across ten topics. Only 72.7\% of generated natural-science citations and 76.6\% of humanities citations were confirmed to exist, with substantial DOI errors and disciplinary differences. \href{https://www.jmir.org/2024/1/e52935/}{source}.
\item
  \textbf{E-EXT-016} (primary-study; ACL abstract and proceedings metadata; questions Q3, Q6): The study evaluates GPT-3.5 and GPT-4 on review scoring and generation using a review-revision dataset. Models can make passable single-option decisions, but fully correct answers are rare and critical feedback on long papers remains weak. \href{https://aclanthology.org/2024.lrec-main.816/}{source}.
\item
  \textbf{E-EXT-017} (systematic-review; publisher abstract, method preview, and results summary; questions Q2, Q3, Q6): Across 25 included empirical studies, effects were more consistently positive for psychological and writing-process outcomes than for textual products. Product gains were clearest for surface accuracy and conventions, while higher-order argumentation and discourse-organization results were mixed. \href{https://www.sciencedirect.com/science/article/pii/S1475158526000986}{source}.
\item
  \textbf{E-EXT-018} (systematic-review; publisher abstract, method preview, and results summary; questions Q2, Q3, Q6): The review reports stronger performance for grammar, coherence, and rubric-based tasks than for higher-order argumentation and creativity, with unresolved validity, fairness, equity, and integrity concerns for consequential assessment. \href{https://www.sciencedirect.com/science/article/pii/S1075293526000292}{source}.
\item
  \textbf{E-EXT-019} (primary-study; official proceedings abstract, full metadata, and linked paper summary; questions Q2, Q4, Q6): Large-scale in-the-wild writing sessions show recurring multi-turn behaviors in which users revise intents, explore text, ask questions, adjust style, or inject content; a small set of interaction patterns explains much of the observed variation. \href{https://aclanthology.org/2025.emnlp-main.852/}{source}.
\item
  \textbf{E-EXT-020} (field-experiment; published article abstract, design summary, and data-availability record; questions Q2, Q4, Q6): Access to a conversational assistant increased average issues resolved per hour by 15 percent, with larger speed and quality gains for less experienced workers and small quality declines among the highest-skilled workers. \href{https://academic.oup.com/qje/article/140/2/889/7990658}{source}.
\item
  \textbf{E-EXT-021} (primary-study; arXiv abstract and metadata; accepted-paper claim not independently checked in proceedings; questions Q2, Q3, Q4): The paper proposes interaction-aware measures of hierarchical suggestion acceptance and knowledge-aware editing distance and reports that they expose behavior and editing cost that output-only metrics miss. \href{https://arxiv.org/abs/2605.23535}{source}.
\item
  \textbf{E-METHOD-004} (governance-framework; official NIST profile page and PDF summary; questions Q5, Q6): NIST's Generative AI Profile extends the AI RMF with risks and actions across the generative-AI lifecycle, including validity, reliability, privacy, transparency, and accountability concerns. \href{https://www.nist.gov/itl/ai-risk-management-framework}{source}.
\item
  \textbf{E-METHOD-005} (reporting-guideline; official PRISMA statement and checklist pages; questions Q6): PRISMA 2020 specifies reporting items for the rationale, methods, study selection, characteristics, risk of bias, synthesis, and limitations of systematic reviews, with flow-diagram templates. \href{https://www.prisma-statement.org/prisma-2020}{source}.
\item
  \textbf{E-METHOD-006} (publication-policy; official policy page; questions Q4, Q5, Q6): ICMJE states that AI tools should not be listed as authors, that humans remain responsible for accuracy, integrity, originality, citations, and plagiarism checks, and that AI use should be disclosed. \href{https://www.icmje.org/recommendations/browse/artificial-intelligence/ai-use-by-authors.html}{source}.
\item
  \textbf{E-SW-003} (software-documentation; official feature documentation; questions Q5): Vale parses markup-aware source formats and can exclude code, URLs, and syntax while applying reusable prose rules. \href{https://vale.sh/features/markup}{source}.
\item
  \textbf{E-SW-004} (software-documentation; official open-source project documentation; questions Q5): LanguageTool exposes an extensible open-source grammar and style engine with rule-development documentation and integrations. \href{https://languagetool.org/dev}{source}.
\item
  \textbf{E-SW-005} (software-documentation; official repository README; questions Q5): LaTeX Lint targets common LaTeX and academic-writing issues and can detect nonstandard mathematical notation and text patterns in TeX and Markdown. \href{https://github.com/HirokiHamaguchi/latexlint}{source}.
\item
  \textbf{E-SW-006} (software-documentation; official filter and AST documentation; questions Q5, Q6): Pandoc exposes a typed document abstract syntax tree to filters, allowing transformations over structural elements rather than unstructured source strings. \href{https://pandoc.org/filters.html}{source}.
\item
  \textbf{E-SW-007} (software-documentation; official API documentation; questions Q5, Q6): pylatexenc provides Python APIs for parsing LaTeX into nodes and for text conversion, including customizable macro and environment handling. \href{https://pylatexenc.readthedocs.io/en/latest/}{source}.
\item
  \textbf{E-SW-008} (software-documentation; official CTAN package record and documentation summary; questions Q5): latexdiff marks textual additions and deletions in two LaTeX sources and can generate a typeset visual comparison. \href{https://www.ctan.org/pkg/latexdiff}{source}.
\item
  \textbf{E-SW-009} (software-documentation; official documentation introduction; questions Q5): GROBID extracts structured bibliographic and full-text information from scholarly PDFs and exposes service interfaces for document processing. \href{https://grobid.readthedocs.io/en/latest/Introduction/}{source}.
\item
  \textbf{E-SW-010} (software-documentation; official repository documentation; questions Q3, Q5): textstat implements established formula-based readability, reading-time, syllable, sentence, and lexical statistics behind a Python API. \href{https://github.com/textstat/textstat}{source}.
\item
  \textbf{E-SW-011} (software-documentation; official framework documentation; questions Q3, Q5, Q6): Inspect AI provides structured datasets, solvers, scorers, logs, and model-provider interfaces for reproducible model evaluations. \href{https://inspect.aisi.org.uk/}{source}.
\item
  \textbf{E-SW-012} (software-documentation; official repository documentation; questions Q5, Q6): llama.cpp supports local inference for multiple model families and hardware targets through a portable runtime and server interface. \href{https://github.com/ggml-org/llama.cpp}{source}.
\item
  \textbf{E-SW-013} (software-documentation; official successor repository, documentation, and release record; questions Q5, Q6): LTeX+ LS is the maintained successor to the archived LTeX LS repository. It provides offline LanguageTool-backed diagnostics for LaTeX, Markdown, BibTeX, and other markup through LSP and command-line interfaces. \href{https://github.com/ltex-plus/ltex-ls-plus}{source}.
\item
  \textbf{E-SW-014} (software-documentation; official documentation and DOM API overview; questions Q5): plasTeX is a Python LaTeX processor that expands macros into an XML-like document object model and supports custom packages and renderers. \href{https://plastex.github.io/plastex/}{source}.
\item
  \textbf{E-SW-015} (software-documentation; official repository README and stated limitations; questions Q5): tree-sitter-latex supplies an incremental concrete-syntax grammar focused on language-server constructs and explicitly uses best-effort parsing because general TeX macro and catcode behavior cannot be represented by the grammar. \href{https://github.com/latex-lsp/tree-sitter-latex}{source}.
\item
  \textbf{E-SW-016} (software-documentation; official NIST-hosted project manual and repository documentation; questions Q5): LaTeXML converts TeX and LaTeX documents into structured XML, HTML, and MathML with bindings for supported packages. \href{https://dlmf.nist.gov/LaTeXML/}{source}.
\item
  \textbf{E-METHOD-007} (public-database-export; full local export and hash sidecar; questions Q6): The collector preserves the compressed public bibliography hash and applies a declared title filter without treating the resulting records as included studies. \href{https://aclanthology.org/anthology.bib.gz}{source}.
\item
  \textbf{E-METHOD-008} (public-database-export; full local export and hash sidecar; questions Q6): The public NBER search endpoint was queried with a recorded workplace-generative-AI query; the current matching-record count and raw response hash are preserved in the provenance sidecar rather than frozen into the claim. \href{https://www.nber.org/api/v1/search}{source}.
\item
  \textbf{E-METHOD-009} (public-database-export; full local export and hash sidecar; questions Q4, Q6): The official proceedings search returns the diversity study and records the query, result URL, and raw HTML hash. \href{https://proceedings.iclr.cc/papers/search?q=writing+language+models}{source}.
\end{itemize}

\hypertarget{limitations}{%
\subsection{Limitations}\label{limitations}}

The current evidence rows are a curated, provisional synthesis and not a completed systematic review. Abstract-level records cannot support methods-level claims. Public API coverage is not equivalent to EconLit, RePEc, SSRN, ACM, IEEE, Web of Science, or Scopus coverage. The internal case record is valuable for product requirements but is not an external efficacy study. These limits are why the proposal remains gated.
\end{document}
